\section{跨论文的系统性比较与讨论}
\label{sec:synthesis}

六篇论文共同面对的问题，是如何把持续增长的历史转化为后续任务能够维护、找到并正确使用的记忆。但它们干预的环节不同：有的改变信息组织，有的改变访问范围，有的学习证据使用经验，还有的约束记忆操作的执行。因而，综合分析需要同时回答“改进了什么环节”“依赖什么条件”和“仍会在哪一步失败”，再讨论不同机制能否互补。

\subsection{统一比较框架：研究对象、机制与边界}

表\ref{tab:memory-comparison}先对齐研究对象和机制，不将六项工作视为同一个榜单上的六个可直接替换算法。表中列出各论文的主要着力点，并不表示它只具有这一项能力。

\begin{table}[htbp]
    \centering
    \caption{六篇论文的研究定位与比较边界}
    \label{tab:memory-comparison}
    \small
    \begin{tabularx}{\linewidth}{@{}p{0.16\linewidth}XXX@{}}
        \toprule
        工作 & 研究对象与表示 & 主要着力点 & 比较时的边界 \\
        \midrule
        StructMem\upcite{xu-etal-2026-structmem} & 长对话中的事件条目与跨事件综合记忆 & 保留事件内绑定，连接跨事件信息 & 综合阶段的事件重建不等于问答阶段必然采用相同扩展 \\
        Hindsight\upcite{latimer-etal-2026-hindsight} & 区分事实、经历、观察与观点的时序实体记忆 & 写入、融合检索与利用记忆推理 & 端到端系统同时包含组织、检索和推理因素 \\
        Semantic XPath\upcite{liu-etal-2026-semantic} & 树状记忆及其子树与节点 & 将结构范围约束与语义匹配结合 & 依赖可用的层级结构及正确的查询生成 \\
        HyMEM\upcite{zhu-etal-2026-hybrid} & 电脑操作中的离散经验和连续多模态表示 & 混合表示、经验更新和工作记忆刷新 & 网页交互、视觉输入与训练条件不同于纯文本问答 \\
        MetaMem\upcite{xin-etal-2026-metamem} & 事实之外的知识利用经验 & 通过反馈反思优化可复用的符号经验 & 需要区分优化阶段反馈与测试时可用信息 \\
        Text2Mem\upcite{wang-etal-2026-text2mem} & 记忆操作的结构化语言与执行状态 & 明确操作语义，验证规划和执行 & 操作正确性不能直接换算成长对话回答准确率 \\
        \bottomrule
    \end{tabularx}
\end{table}

上述定位分别依据 StructMem 第 2--3 页图 2 与公式 1--6、Hindsight 第 2--5 页图 1、Semantic XPath 第 2--5 页第 3 节、HyMEM 第 3--5 页图 1、MetaMem 第 3--5 页图 2，以及 Text2Mem 第 3--5 页图 2 和表 1。此处页码均按仓库 PDF 从第 1 页起计。

\subsection{机制之间的联系：互补发生在什么条件下}

\textbf{组织结构与访问范围。} StructMem 强调同一事件内的事实与关系绑定，以及跨事件的综合；Semantic XPath 强调查询应落在哪个结构范围。二者可以用来区分“存储时丢失关系”和“读取时选错范围”这两类问题：即使存储正确，也可能检索错会话或错节点；即使检索范围正确，也可能因为原始抽取已遗漏关系而无法回答。这里的互补是分析思路，不能据此认定把 XPath 接到事件记忆上就一定有效；还需定义结构映射、证据来源与额外查询成本。

\textbf{取回证据与使用证据。} Hindsight 的 recall 与 reflect 分别涉及检索和利用记忆推理；MetaMem 则把可复用的知识利用经验作为优化对象。它们共同提示：检索命中后仍可能出现重复计数、时间关系误读或不当综合。分析回答错误时，需要区分“证据没有进入上下文”和“证据已进入但被用错”。改善前者的检索策略与改善后者的利用经验可能互补，但效果必须分别验证，不能只观察最终分数便认定瓶颈已经解决。

\textbf{更新策略与操作语义。} HyMEM 在具体任务中决定如何合并、替换和刷新经验；Text2Mem 则强调操作应具有明确的目标、参数与可检查的执行结果。前者帮助讨论“何时更新、保留什么”，后者帮助讨论“实际执行了什么、状态是否符合意图”。但二者的表示与接口并不天然一致，尤其不能把多模态连续记忆直接当作参考 SQL 后端的一条文本记录。

\textbf{事实、观点与利用经验。} Hindsight 对事实与观点的区分，以及 MetaMem 对知识利用经验的建模，说明记忆条目的角色需要明确。世界事实、对主体的推断、操作经验和回答策略不宜无差别地作为同一类证据。它们分别在什么条件下可用于回答、何时需要更新或撤销，是比较系统设计时应追问的问题，而非仅看数据结构叫“图”还是“树”。

\subsection{共同取舍：结构化获得了什么，又付出了什么}

\textbf{关联完整性与抽取误差。} 事件、实体关系和层级有助于保留上下文，但结构本身也由抽取或整理得到。实体合并错误、时间绑定错误及摘要遗漏会进入后续检索和推理。因此应同时检查结构是否找到了来源，以及来源中的关键事实是否仍被保留，不能用“有来源 ID”替代内容正确性。

\textbf{检索覆盖与上下文开销。} 多路召回、事件补齐和更大范围的结构访问可能增加有效证据，也会引入冗余和无关信息。评价局部改动时，应在相同回答上下文预算下观察证据覆盖、回答质量及成本，避免把“给模型更多文本”的收益全部归因于结构设计。

\textbf{动态适应与错误固化。} 综合记忆、观点更新和反思经验能够适应后续任务，也可能保留错误概括或过拟合开发样本。讨论应追溯更新依赖的证据和反馈，区分系统运行中自然获得的信息与离线优化时使用的标准答案；后者不能出现在独立测试的记忆构建或检索条件中。

\textbf{表达能力与验证范围。} 明确操作 schema 能减少歧义、支持执行检查，但可解析不代表目标选择正确，可执行也不代表执行后满足用户意图。同样，能组织复杂多模态经验，不代表在更简单的文本任务上具有成本优势。需要分别报告表示、访问、推理和执行各环节的约束。

以上是用于比较的分析维度，不是声称所有论文都已经验证了这些取舍的同方向收益。具体优缺点需回到各篇论文的实验、消融和局限，避免把设计动机当成已证明的结论。

\subsection{怎样比较实验结果，才支持方法判断}

六篇论文的任务覆盖长对话问答、结构化任务状态、电脑操作和记忆指令执行。回答准确率、网页任务成功率、操作执行率和状态匹配指标测量的对象不同，不能汇总为一个“总体性能”分数。即使都报告 LoCoMo，模型、数据处理、题型范围、检索与回答预算、评分提示词等仍可能不同。

比较时应记录任务与数据划分、模型角色、是否训练或使用优化反馈、检索范围和预算、评价指标、主要成本六类条件。原论文的跨系统结果说明其各自实验条件下的表现；要判断某个结构或机制的作用，优先看该论文内部控制变量的消融。若本组复现，则所有对照共用明确的数据划分和评价口径，并说明与原论文配置的差异。

成本也要按机制分开：写入与综合开销、检索和回答开销、优化与裁判开销不是同一项。无需训练基座模型，并不等于无需多次采样与反馈；一次构建的成本，也不能直接与一次查询的成本相比。

\subsection{用共同案例检查六种研究视角}

可围绕一个明确标注为自拟的讨论案例展开：历史记录中，某人先计划周五到 A 地，后来改为周六到 B 地；随后提到从 B 地继续出发。查询要求回答最终出行时间和后续路线。

这个案例首先涉及事件内人物、时间和地点是否绑定正确，以及跨事件关系是否能被找回；其次涉及检索是否选对会话、时间和状态版本；还涉及已取回新旧计划后，是否能正确处理变更，而非把两个版本拼成矛盾答案。若用户再要求暂时不提及这次出行，则还要明确这是降低检索优先级、临时隐藏还是删除，操作范围和期限分别是什么。

讨论时应由各论文负责人说明：其方法直接处理哪一部分、依赖什么输入、哪一部分并未给出保证，并指出原文依据。这个案例用于检验解释是否一致，不是论文原例，也不是已运行的六系统对照实验；尤其不能把纯文本例子直接视为验证了 HyMEM 的多模态机制。

\subsection{从比较形成课设选择与团队结论}

综合分析最终应落实为有依据的选择：本组要优先解决信息组织、检索范围、证据利用还是操作执行中的哪一个问题；为何选择某项工作作为主复现对象；哪些机制只作为参考，哪些因任务或资源不匹配而暂不纳入。选择某篇作为主线并不意味着其他五篇只作背景罗列，还应说明它们如何帮助定义错误类型、评价指标或可选改进。

此部分等待实际团队讨论后补充：保留有原文或实验依据的共同结论；对尚有分歧的问题写明分歧点与需要核查的证据。若提出机制组合，应分别列出结构适配、信息边界和对照实验，而非直接宣称已提出有效的新方法。最终报告与 PPT 应呈现这些经过讨论的取舍，并在成员分工中记录实际参与的比较、核对和整合工作。
